% Полуторный интервал и красная строка 1.25 см
\usepackage{setspace}
\onehalfspacing
\usepackage{indentfirst}
\setlength{\parindent}{1.25cm}

% Нумерация страниц внизу по центру
\usepackage{fancyhdr}
\pagestyle{fancy}
\fancyhf{}
\fancyfoot[C]{\thepage}
\renewcommand{\headrulewidth}{0pt}
\fancypagestyle{plain}{
  \fancyhf{}
  \fancyfoot[C]{\thepage}
  \renewcommand{\headrulewidth}{0pt}
}

% Заголовки глав: "ГЛАВА 1. НАЗВАНИЕ" (как в образце)
\usepackage{titlesec}
\titleformat{\chapter}[block]
  {\normalfont\bfseries\fontsize{14}{16.8}\selectfont\MakeUppercase}
  {ГЛАВА \thechapter.}{0.5em}{}
\titlespacing*{\chapter}{1.25cm}{0pt}{12pt}

% Заголовки параграфов: "1.1 Название" (без точки после номера)
\titleformat{\section}[block]
  {\normalfont\bfseries\fontsize{14}{16.8}\selectfont}
  {\thesection}{0.5em}{}
\titlespacing*{\section}{1.25cm}{12pt}{6pt}

\titleformat{\subsection}[block]
  {\normalfont\bfseries\fontsize{14}{16.8}\selectfont}
  {\thesubsection}{0.5em}{}
\titlespacing*{\subsection}{1.25cm}{12pt}{6pt}

% Подписи к рисункам (снизу, по центру) и таблицам (сверху, слева)
\usepackage{caption}
\captionsetup[figure]{
  justification=centering,
  labelsep=endash,
  font={rm, normalsize}
}
\captionsetup[table]{
  justification=raggedright,
  singlelinecheck=false,
  labelsep=endash,
  font={rm, normalsize}
}

% Надпись "Продолжение таблицы X.Y" при переносе длинных таблиц (longtable)
\usepackage{etoolbox}
\AtBeginEnvironment{longtable}{\small}